% GENERATED FILE. Edit canonical lessons, then run npm run generate:books.
% canonical-source-hash: fnv1a64:e06391c5cf61ce5a
% canonical-lessons: FA-C524-kargardan, FA-C524-filmname, FA-C524-filmbardar, FA-C524-naqsh-bazi-kardan, FA-C524-honarpishe

\chapter{Making Films}
\label{ch:fa-making-films}
\hlchaptermodality{524}

\begin{chapteropening}
\textbf{By the end of this chapter:} \emph{I can talk about how films are made.}

Complete the last lesson of chapter 524: I can talk about how films are made.
\end{chapteropening}

\section[\texorpdfstring{\fa{کارگردان} (kârgardân)}{کارگردان (kârgardân)}]{\fa{کارگردان} (kârgardân) — a film director}

\label{lesson:FA-C524-kargardan}



\begin{quote}
Before the new one: say the Persian for a spectator, then the Persian for a role.
\end{quote}

\subsection*{You'll want to know: \fa{کارگردان}}
\textbf{\fa{کارگردان}} (\emph{kârgardân}) — \textquotedblleft{}a film director\textquotedblright{}.

\begin{quote}
\fa{کارگردان} \fa{این} \fa{فیلم} \fa{جایزه} \fa{گرفت}.
\end{quote}

\begin{quote}
\emph{kârgardân-e in film jâyeze gereft.}
\end{quote}

\begin{quote}
The director of this film won a prize.
\end{quote}

\emph{kâr} 'work' + \emph{gardân} 'turning' (stem of \emph{gardândan}): the one who keeps the work turning.

\begin{cousinweb}
\emph{kârgardân} is built on native Persian \emph{kâr}, 'work', which EtymDB traces to Middle Persian (EtymDB), and \emph{gardân}, the stem of \emph{gardândan}, 'to turn, to run'.
\end{cousinweb}

\begin{grammarlens}[title={What you've built}]
One more word for telling, explaining and describing at B1.
\end{grammarlens}

\subsection*{Your turn}
\begin{itemize}
\raggedright
  \item \emph{Say it:} \emph{kârgardân}
  \item \emph{Say it:} \emph{kârgardân}, once more
  \item \emph{Say it:} a sentence of your own with \emph{kârgardân}
  \item \emph{Recall:} say the Persian for a spectator, then the Persian for a role, then say \emph{kârgardân} again
\end{itemize}

\begin{tcolorbox}[breakable,colback=teal!4,colframe=teal!35!black,title={Before you move on}]
What does \fa{کارگردان} mean? (\textquotedblleft{}A film director\textquotedblright{}.) Say it once more.
\end{tcolorbox}
\section[\texorpdfstring{\fa{فیلمنامه} (filmnâme)}{فیلمنامه (filmnâme)}]{\fa{فیلمنامه} (filmnâme) — a screenplay, a film script}

\label{lesson:FA-C524-filmname}



\begin{quote}
Before the new one: say the Persian for a role, then the Persian for a film director.
\end{quote}

\subsection*{You'll want to know: \fa{فیلمنامه}}
\textbf{\fa{فیلمنامه}} (\emph{filmnâme}) — \textquotedblleft{}a screenplay, a film script\textquotedblright{}.

\begin{quote}
\fa{او} \fa{فیلمنامه} \fa{را} \fa{در} \fa{یک} \fa{ماه} \fa{نوشت}.
\end{quote}

\begin{quote}
\emph{u filmnâme râ dar yek mâh nevesht.}
\end{quote}

\begin{quote}
He wrote the screenplay in a month.
\end{quote}

\emph{film} + \emph{nâme} 'writing'. A screenwriter is \emph{filmnâme-nevis}.

\begin{cousinweb}
\emph{filmnâme} is built in Persian on \emph{film}, a loan usually traced to English \emph{film}, and native \emph{nâme}, 'letter, writing'.
\end{cousinweb}

\begin{grammarlens}[title={What you've built}]
One more word for telling, explaining and describing at B1.
\end{grammarlens}

\subsection*{Your turn}
\begin{itemize}
\raggedright
  \item \emph{Say it:} \emph{filmnâme}
  \item \emph{Say it:} \emph{filmnâme}, once more
  \item \emph{Say it:} a sentence of your own with \emph{filmnâme}
  \item \emph{Recall:} say the Persian for a role, then the Persian for a film director, then say \emph{filmnâme} again
\end{itemize}

\begin{tcolorbox}[breakable,colback=teal!4,colframe=teal!35!black,title={Before you move on}]
What does \fa{فیلمنامه} mean? (\textquotedblleft{}A screenplay, a film script\textquotedblright{}.) Say it once more.
\end{tcolorbox}
\section[\texorpdfstring{\fa{فیلمبردار} (filmbardâr)}{فیلمبردار (filmbardâr)}]{\fa{فیلمبردار} (filmbardâr) — a cameraman, a camera operator}

\label{lesson:FA-C524-filmbardar}



\begin{quote}
Before the new one: say the Persian for a film director, then the Persian for a screenplay.
\end{quote}

\subsection*{You'll want to know: \fa{فیلمبردار}}
\textbf{\fa{فیلمبردار}} (\emph{filmbardâr}) — \textquotedblleft{}a cameraman, a camera operator\textquotedblright{}.

\begin{quote}
\fa{فیلمبردار} \fa{از} \fa{کوه} \fa{فیلم} \fa{گرفت}.
\end{quote}

\begin{quote}
\emph{filmbardâr az kuh film gereft.}
\end{quote}

\begin{quote}
The camera operator filmed the mountain.
\end{quote}

\emph{film} + \emph{bardâr} 'taking' (stem of \emph{bardâshtan}).

\begin{cousinweb}
\emph{filmbardâr} is built in Persian on \emph{film}, a loan usually traced to English, and \emph{bardâr}, the stem of native \emph{bardâshtan}, 'to take up'.
\end{cousinweb}

\begin{grammarlens}[title={What you've built}]
One more word for telling, explaining and describing at B1.
\end{grammarlens}

\subsection*{Your turn}
\begin{itemize}
\raggedright
  \item \emph{Say it:} \emph{filmbardâr}
  \item \emph{Say it:} \emph{filmbardâr}, once more
  \item \emph{Say it:} a sentence of your own with \emph{filmbardâr}
  \item \emph{Recall:} say the Persian for a film director, then the Persian for a screenplay, then say \emph{filmbardâr} again
\end{itemize}

\begin{tcolorbox}[breakable,colback=teal!4,colframe=teal!35!black,title={Before you move on}]
What does \fa{فیلمبردار} mean? (\textquotedblleft{}A cameraman, a camera operator\textquotedblright{}.) Say it once more.
\end{tcolorbox}
\section[\texorpdfstring{\fa{نقش} \fa{بازی} \fa{کردن} (naqsh bâzi kardan)}{نقش بازی کردن (naqsh bâzi kardan)}]{\fa{نقش} \fa{بازی} \fa{کردن} (naqsh bâzi kardan) — to play a role, to act a part}

\label{lesson:FA-C524-naqsh-bazi-kardan}



\begin{quote}
Before the new one: say the Persian for a screenplay, then the Persian for a cameraman.
\end{quote}

\subsection*{You'll want to know: \fa{نقش} \fa{بازی} \fa{کردن}}
\textbf{\fa{نقش} \fa{بازی} \fa{کردن}} (\emph{naqsh bâzi kardan}) — \textquotedblleft{}to play a role, to act a part\textquotedblright{}.

\begin{quote}
\fa{او} \fa{در} \fa{این} \fa{فیلم} \fa{نقش} \fa{یک} \fa{پزشک} \fa{را} \fa{بازی} \fa{میکند}.
\end{quote}

\begin{quote}
\emph{u dar in film naqsh-e yek pezeshk râ bâzi mikonad.}
\end{quote}

\begin{quote}
She plays a doctor in this film.
\end{quote}

\emph{naqsh-e kasi râ bâzi kardan} — to play someone's part.

\begin{cousinweb}
\emph{naqsh bâzi kardan} is built in Persian on \emph{naqsh}, 'role', borrowed from Arabic \emph{naqsh}, and the native phrase \emph{bâzi kardan}, 'to play'.
\end{cousinweb}

\begin{grammarlens}[title={What you've built}]
One more word for telling, explaining and describing at B1.
\end{grammarlens}

\subsection*{Your turn}
\begin{itemize}
\raggedright
  \item \emph{Say it:} \emph{naqsh bâzi kardan}
  \item \emph{Say it:} \emph{naqsh bâzi kardan}, once more
  \item \emph{Say it:} a sentence of your own with \emph{naqsh bâzi kardan}
  \item \emph{Recall:} say the Persian for a screenplay, then the Persian for a cameraman, then say \emph{naqsh bâzi kardan} again
\end{itemize}

\begin{tcolorbox}[breakable,colback=teal!4,colframe=teal!35!black,title={Before you move on}]
What does \fa{نقش} \fa{بازی} \fa{کردن} mean? (\textquotedblleft{}To play a role, to act a part\textquotedblright{}.) Say it once more.
\end{tcolorbox}
\section[\texorpdfstring{\fa{هنرپیشه} (honarpishe)}{هنرپیشه (honarpishe)}]{\fa{هنرپیشه} (honarpishe) — an actor, an actress}

\label{lesson:FA-C524-honarpishe}



\begin{quote}
Before the new one: say the Persian for a cameraman, then the Persian for to play a role.
\end{quote}

\subsection*{You'll want to know: \fa{هنرپیشه}}
\textbf{\fa{هنرپیشه}} (\emph{honarpishe}) — \textquotedblleft{}an actor, an actress\textquotedblright{}.

\begin{quote}
\fa{این} \fa{هنرپیشه} \fa{در} \fa{تئاتر} \fa{هم} \fa{بازی} \fa{میکند}.
\end{quote}

\begin{quote}
\emph{in honarpishe dar te'âtr ham bâzi mikonad.}
\end{quote}

\begin{quote}
This actor also works in the theatre.
\end{quote}

\emph{honar} 'art' + \emph{pishe} 'trade'. It means the same as \emph{bâzigar}.

\begin{cousinweb}
\emph{honarpishe} is built on native Persian \emph{honar}, 'art', usually traced to Middle Persian \emph{hunar}, and \emph{pishe}, 'trade, craft': one whose trade is art.
\end{cousinweb}

\begin{grammarlens}[title={What you've built}]
One more word for telling, explaining and describing at B1.
\end{grammarlens}

\subsection*{Your turn}
\begin{itemize}
\raggedright
  \item \emph{Say it:} \emph{honarpishe}
  \item \emph{Say it:} \emph{honarpishe}, once more
  \item \emph{Say it:} a sentence of your own with \emph{honarpishe}
  \item \emph{Recall:} say the Persian for a cameraman, then the Persian for to play a role, then say \emph{honarpishe} again
\end{itemize}

\begin{tcolorbox}[breakable,colback=teal!4,colframe=teal!35!black,title={Before you move on}]
What does \fa{هنرپیشه} mean? (\textquotedblleft{}An actor, an actress\textquotedblright{}.) Say it once more.
\end{tcolorbox}
